% \chapter{Chapter 1}\label{cha:cha1}
Due to its unprecedented sensitivity in the infrared (IR) bands, the James Webb Space Telescope (JWST) has detected several active galactic nuclei (AGNs) at high redshift $(z\sim3-11)$ with much lower bolometric luminosities than those studied previously~\cite{Maiolini_2025}.
Such objects are called `little red dots' (LRDs).
An important property of these objects is their X-ray weakness, as they have been undetected at the limits of the deepest Chandra fields\footnote{Field in which the most sensitive X-ray surveys were conducted.}, regardless of their optical classification and even when stacking signals from multiple sources, with only very few exceptions.
Since the discovery of LRDs in 2024~\cite{matthee2024littlereddotsabundant}, multiple theories to explain the X-ray weakness have been proposed.
There are three main theories: the measurable X-ray spectra are heavily absorbed; LRDs are intrinsically X-ray weak; or the AGNs were misidentified. \newline
The last hypothesis can be ruled out as very unlikely, since the same X-ray weakness has been found independently of the AGN identification method and because there are objects very clearly identified as AGNs that are extremely X-ray weak~\cite{Maiolini_2025}.
Furthermore, LRDs have broad $\text H\alpha$ emission lines, which have been proposed as possible misidentifications, as other objects can produce similar emission lines, including strong outflows driven by star formation or supernovae.
Even though this hypothesis agrees with MIRI observations~\cite{PerezGonzalez_2024}, NIRSpec observations have shown characteristics of LRDs that are hard to explain with this scenario~\cite{Maiolini_2025, Yue_2024}. 
Future studies of emission lines in large LRD samples will give more information about the contribution of outflows. \newline
The following scenarios adopt two different types of AGN\@.
Type 1 AGN have broad emission lines, primarily consisting of $\text H_\alpha$ and $\text H_\beta$, while these emission lines are narrower for LRDs identified as type 2 AGN\@.

\noindent
The heavy-absorption scenario is based on gas column densities within the Compton-thick\footnote{Compton-thick refers to the effect that Compton-scattering dampens the X-Ray spectrum in high density gas-clouds.} regime that absorb the measurable X-ray spectra.
Observing the $\text H_\alpha$ and $\text H_\beta$ emission lines in the presence of such absorption requires extremely large gas densities, that are typical of the broad line region (BLR) with hydrogen column densities of $N_{\text H}=\SI{2e24}{\per\centi\meter\squared}$ and a covering factor approaching $4\pi$,\cite{refId0}.
To achieve such a large covering factor, these clouds are suggested to house low metallicities, which imply low dust content and therefore reduces the radiation pressure feedback of the AGN\@.
As a result, dust-poor gas accumulates around the nuclear region\cite{Ishibashi_2025}.
This would also explain the reduced equivalent widths of $\text{Fe}_{\text{II}}$ and $\text H_\beta$ lines found in LRDs~\cite{Ishibashi_2025}.
A recent study by Comastri et al. (2026)~\cite{refId0} found type 2 AGN with a significant detection of $\sim3\sigma$ and a total stacked exposure time of $\SI{210}{\mega\second}$ in the ultra-hard band (UHB) of the X-ray spectrum ($10-30\ \text{keV}$), but not in lower energy bands.
This supports the hypothesis of heavy-absorption as only highly-energetic X-ray photons can pass through the Compton-thick absorber, while lower energetic photons remain undetected.
The same procedure for type 1 LRDs has not revealed any evidence of a significant signal in about $\SI{140}{\mega\second}$ stacked data, suggesting that they are additionally intrinsically X-ray weak.
Other studies have ruled out heavy-absorption as the only compatible solution, at high levels of obscuration with column densities of $N_{\text H} = 10^{24.5-25}\ \si{\per\centi\meter\squared}$\cite{Sacchi_2025} that are only possible when the obscuration occurs in the same clouds responsible for the broad-line emission~\cite{Maiolini_2025}.
Current objects found with such high column densities of $\gtrsim10^{24}$ are broad absorption line (BAL) quasars, but LRDs do not show any strong BAL features, arguing against the heavy-absorption theory~\cite{Yue_2024}. \newline
The intrinsically weak X-ray emission theory is based on highly accreting black holes (BH) that lead to lower temperatures close to the BH and a weak jet with lower spin, resulting in a softer X-ray spectrum.
Super-Eddington accretion rates lead to stronger winds, that produce a broader BLR\@.
Chandra may only be able to detect AGNs at minimal angles from the jet, as the detection rate of a supermassive black hole (SMBH) within $\SI{10}{\degree} (\sim1.5\%)$ is similar to detecting X-rays by LRDs $(\sim0.6\%)$~\cite{Pacucci_2024}.
Furthermore, narrow-line Seyfert 1 (NLSy1) galaxies are known to have high accretion rates close to the Eddington limit, or even super-Eddington, and show steeper X-ray spectra than normal AGN~\cite{Maiolini_2025}.
This supports the theory by Comastri et al. (2026)~\cite{refId0} that type 1 AGN are likely to be obscured and intrinsically weak, as the similarity with NLSy1 cannot account for the entire X-ray weakness.
A new study by Sacchi et al. (2025)~\cite{Sacchi_2025}, however, found no detection despite a high exposure time of $\SI{400}{\mega\second}$.
The calculated upper limits are deep enough to rule out current super-Eddington accretion models of Pacucci et al.~\cite{Pacucci_2024} and Inayoshi et al.~\cite{inayoshi2025weaknessxraysvariabilityhighredshift}, stating that the X-ray weakness is only explainable with heavy-obscuration.

A major limitation of studies of the X-ray weakness of LRDs is that only a few photons have been detected and that analyses rely on stacked data.
In a stacked dataset, some LRDs could be intrinsically weak, while others may be heavily obscured.
This leads to the current consensus that a combination of both leading hypotheses will most likely be true for type 1 AGN\@. 
Evidence of $\sim3\sigma$ significance has been found that heavy-obscuration is causing the X-ray weakness in type 2 AGNs~\cite{refId0}.
Intrinsic X-ray weakness via super-Eddington accretion will never be the sole solution, because the upper-limit calculated by Sacchi et al.~\cite{Sacchi_2025} may rule out current theories.
It remains unclear within which framework the type 1 AGNs will fit.
A decisive future observation would be the detection of type 1 AGNs in the UHB, likely ruling out super-Eddington accretion and confirming the heavy-obscuration scenario.
This detection might be published with data taken from next-generation observatories like the Advanced X-ray Imaging Satellite (AXIS) hosting detectors with a bandpass of $0.3-10\ \si{\kilo\eV}$ and a spatial resolution of $\sim\SI{1.5}{\arcsecond}$ across the entire field of view~\cite{Reynolds_2023}.
